\part{Risiken, Empfehlungen und Abschluss}
\chapterimage{Bilder_und_Co/safety_sign_with_computer_code_and_a_airplane_in_the_background_1958542794.png} % Chapter heading image
%----------------------------------------------------------------------------------------
\chapter{Sicherheitslücken und Risiken}\label{sec:luecken-risiken}
Auf Basis der Analyse werden Lücken und Risiken konsolidiert und priorisiert. Die Darstellung orientiert sich an Anforderungsfeldern und benennt jeweils Ursache, Auswirkung und mögliche Gegenmaßnahmen auf hoher Ebene. % TODO: Liste/Matrix mit Priorität
  \section{Konsolidierte Befunde}
  \section{Priorisierung}










\chapter{Empfehlungen zur Anwendung}\label{sec:empfehlungen}
Die Empfehlungen leiten pragmatische Maßnahmen für kurz-, mittel- und langfristige Zeithorizonte ab. Im Vordergrund stehen operationalisierbare Schritte zur Stärkung von Evidenzketten, Traceability und Prozessabsicherung. % TODO: 3–3–3 Maßnahmen (kurz)
  \section{Kurzfristige Maßnahmen}
  \section{Mittelfristige Entwicklungspfade}
  \section{Langfristige Standardisierung}











\chapter{Fazit und Ausblick}\label{sec:fazit}
Das Fazit bündelt die Antwort auf die Leitfrage im Rahmen der getroffenen Einordnung und verweist auf den weiteren Entwicklungsbedarf. Der Ausblick skizziert potenzielle Vertiefungen und den erwartbaren Reifegewinn durch Standardisierung und Best Practices. % TODO: 5–7 Sätze, konsistent mit Abstract/Ziel
